\documentclass[a4paper]{article}

\usepackage{graphicx}
\usepackage{amsmath,amssymb}
\usepackage{minted}
\usepackage{multirow}
\usepackage{float}
\usepackage{todonotes}
\usepackage{forest}
\usepackage{xstring}
\usepackage{xcolor}
\usepackage[most]{tcolorbox}
\usepackage{xparse}
\usepackage{natbib}

\usetikzlibrary{graphs,graphdrawing}
\usegdlibrary{layered}

% for external graphviz
\input{/home/donkarlo/Dropbox/repo/nd_ai_project/src/nd_ai/knoweldge_representation/language/formal/computer/programming/tex/latex/code_clips/external_graphviz.tex}

% for tags
\input{/home/donkarlo/Dropbox/repo/nd_ai_project/src/nd_ai/knoweldge_representation/language/formal/computer/programming/tex/latex/parts/control_sequence/kind/semtag/semtag.tex}

% for links
\input{/home/donkarlo/Dropbox/repo/nd_ai_project/src/nd_ai/knoweldge_representation/language/formal/computer/programming/tex/latex/code_clips/seehere.tex}

\title{Proprioceptive sensors}
\date{}

\begin{document}
    \maketitle
    
    
    \section{Proprioception}
        
        \clickurlfooter{https://en.wikipedia.org/wiki/Proprioception} defines is the sense of self-movement, force, and body position.
        
        \cite{tuthill-2018-proprioception} defines proprioception as the sense that lets the body monitor position, movement, and force, even when we are not consciously aware of it.
        The article explains that this sense is essential for coordinated action, and shows that losing proprioception can severely disrupt even basic movement.
        It also reviews the sensory receptors and neural pathways that allow the nervous system to build an internal representation of the body.
        \bibliographystyle{plainnat}
        \bibliography{/home/donkarlo/Dropbox/repo/nd_shared_research/bibliography/bibliography.bib}
\end{document}